\documentclass[12pt]{article}
\usepackage{amsmath, amssymb, amsfonts, geometry, setspace, tcolorbox}
\geometry{margin=1in}
\setstretch{1.2}
\setlength{\parindent}{0pt}


\title{Optimization of Multi-Variable Functions\\[7pt]
\large ECON 320 - Introduction to Mathematical Economics}
\author{Muhebullah Karimzada}
\date{Fall 2025}

\begin{document}
\maketitle

\section*{1. Introduction}

Many economic problems require maximizing or minimizing a function of several variables. 
For example, a firm maximizes profit depending on capital and labor, or a consumer maximizes utility depending on two goods. 
This is the domain of \textbf{multivariable optimization}.

\[
\text{Maximize or minimize } f(x_1, x_2, \ldots, x_n)
\]
possibly \textbf{subject to constraints} such as a budget or production function.



\section*{2. Unconstrained Optimization in Two Variables}

Let
\[
z = f(x, y)
\]
represent a smooth surface.

\subsection*{2.1 First-Order (Necessary) Conditions}

At an interior extremum, small changes in \(x\) or \(y\) leave \(z\) unchanged:

\[
f_x = \frac{\partial f}{\partial x} = 0, \quad f_y = \frac{\partial f}{\partial y} = 0
\]

The points \((x^*, y^*)\) satisfying these equations are called \textit{stationary points}.



\subsection*{2.2 Second-Order (Sufficient) Conditions}

Let
\[
f_{xx} = \frac{\partial^2 f}{\partial x^2}, \quad f_{yy} = \frac{\partial^2 f}{\partial y^2}, \quad f_{xy} = \frac{\partial^2 f}{\partial x \partial y}
\]
and define the \textbf{Hessian determinant}
\[
D = 
\begin{vmatrix}
f_{xx} & f_{xy} \\
f_{yx} & f_{yy}
\end{vmatrix}
= f_{xx}f_{yy} - (f_{xy})^2
\]

\begin{center}
\begin{tabular}{c|c}
\textbf{Condition} & \textbf{Nature of Stationary Point} \\ \hline
$D>0$, $f_{xx}<0$ & Local Maximum \\
$D>0$, $f_{xx}>0$ & Local Minimum \\
$D<0$ & Saddle Point \\
$D=0$ & Inconclusive
\end{tabular}
\end{center}



\subsection*{Example 1: Profit Maximization}

\[
\pi(x,y) = 100x + 80y - 10x^2 - 4y^2 - 4xy
\]

\[
\pi_x = 100 - 20x - 4y = 0, \qquad \pi_y = 80 - 8y - 4x = 0
\]

Solve:
\[
\begin{cases}
20x + 4y = 100 \\
4x + 8y = 80
\end{cases}
\Rightarrow x=4, \; y=8
\]

Second derivatives:
\[
f_{xx}=-20, \; f_{yy}=-8, \; f_{xy}=-4
\]
\[
D = (-20)(-8) - (-4)^2 = 144 > 0, \; f_{xx}<0
\Rightarrow \text{Local Maximum.}
\]



\section*{3. Generalization to \(n\) Variables}

For \(f(x_1, \ldots, x_n)\):

\[
\frac{\partial f}{\partial x_i} = 0, \quad i = 1,2,\ldots,n
\]

The second-order conditions use the \textbf{Hessian matrix}:
\[
H =
\begin{bmatrix}
f_{x_1x_1} & \cdots & f_{x_1x_n} \\
\vdots & \ddots & \vdots \\
f_{x_nx_1} & \cdots & f_{x_nx_n}
\end{bmatrix}
\]

To determine the nature of the stationary point, evaluate the \textbf{leading principal minors} \(D_1, D_2, \ldots, D_n\):

\begin{center}
\begin{tabular}{c|l}
\textbf{Condition} & \textbf{Type of Extremum} \\ \hline
Alternating signs ($D_1<0, D_2>0, D_3<0,\ldots$) & Maximum \\
All positive ($D_1>0, D_2>0, D_3>0,\ldots$) & Minimum
\end{tabular}
\end{center}



\section*{4. Constrained Optimization: The Lagrange Method}

Most economic problems involve constraints.

\[
\max_{x,y} f(x,y) \quad \text{subject to} \quad g(x,y) = k
\]

\subsection*{4.1 Lagrangian Function}
\[
\mathcal{L}(x,y,\lambda) = f(x,y) + \lambda(k - g(x,y))
\]

\subsection*{4.2 First-Order Conditions (FOCs)}
\[
\begin{cases}
f_x - \lambda g_x = 0 \\
f_y - \lambda g_y = 0 \\
g(x,y) = k
\end{cases}
\]

At the optimum,
\[
\frac{f_x}{g_x} = \frac{f_y}{g_y} = \lambda
\]
so the \textbf{marginal rates of substitution} are equal to the ratio of the constraint effects.



\subsection*{Example 2: Utility Maximization under a Budget Constraint}

\[
U(x,y) = x^{0.5}y^{0.5}, \quad p_xx + p_yy = M
\]

Lagrangian:
\[
\mathcal{L} = x^{0.5}y^{0.5} + \lambda(M - p_xx - p_yy)
\]

FOCs:
\[
\frac{1}{2}x^{-0.5}y^{0.5} = \lambda p_x, \quad
\frac{1}{2}x^{0.5}y^{-0.5} = \lambda p_y
\]

Dividing:
\[
\frac{y}{x} = \frac{p_x}{p_y} \Rightarrow y = \frac{p_x}{p_y}x
\]

Substitute in constraint:
\[
p_xx + p_y\left(\frac{p_x}{p_y}x\right) = M \Rightarrow 2p_xx = M \Rightarrow x^* = \frac{M}{2p_x}, \quad y^* = \frac{M}{2p_y}
\]

\textbf{Interpretation:} The consumer chooses a point where MRS = price ratio.


\section*{5. The Economic Meaning of the Lagrange Multiplier}

The multiplier $\lambda$ represents the \textbf{shadow price} or \textbf{marginal value of relaxing the constraint}.

\[
\frac{\partial f^*}{\partial k} = \lambda
\]

In utility maximization, $\lambda$ equals the marginal utility of income; in production, it represents the marginal output per additional unit of input cost.



\section*{6. Multiple Constraints}

If there are $m$ constraints:
\[
g_i(x_1,\ldots,x_n) = k_i, \quad i=1,\ldots,m
\]
then
\[
\mathcal{L} = f(x_1,\ldots,x_n) + \sum_{i=1}^m \lambda_i(k_i - g_i)
\]

First-order conditions:
\[
f_{x_j} = \sum_{i=1}^m \lambda_i g_{i,x_j}, \quad j = 1,\ldots,n
\]



\section*{7. Second-Order Conditions for Constrained Problems}

For a single constraint, the test uses the \textbf{bordered Hessian:}
\[
H_b =
\begin{vmatrix}
0 & g_x & g_y \\
g_x & f_{xx} & f_{xy} \\
g_y & f_{yx} & f_{yy}
\end{vmatrix}
\]

\begin{center}
\begin{tabular}{c|c}
\textbf{Sign of } $H_b$ & \textbf{Type of Extremum} \\ \hline
$(-1)^{m}det(H_b) < 0$ & Maximum \\
$(-1)^{m}det(H_b)> 0$ & Minimum
\end{tabular}
\end{center}

\textbf{Why the $(-1)^{m}$ factor?}

Because each equality constraint adds one border row and column to the Hessian.
Each border introduces a sign flip in the determinant pattern.

Hence, the test must account for how many constraints 
$m$ we have added.   That's exactly what the factor $(-1)^{m}$ does,  it adjusts the sign convention.

\section*{8. Example 3: Cost Minimization with Cobb–Douglas Technology}

\[
Q = K^{0.5}L^{0.5}, \quad \text{Minimize } C = wL + rK, \quad Q = \bar{Q}
\]

Lagrangian:
\[
\mathcal{L} = wL + rK + \lambda(\bar{Q} - K^{0.5}L^{0.5})
\]

FOCs:
\[
\begin{cases}
w - \frac{\lambda}{2}K^{0.5}L^{-0.5} = 0 \\
r - \frac{\lambda}{2}K^{-0.5}L^{0.5} = 0
\end{cases}
\]

Dividing:
\[
\frac{w}{r} = \frac{K}{L} \Rightarrow \frac{L}{K} = \frac{r}{w}
\]

Substitute into constraint:
\[
K^* = \bar{Q}\sqrt{\frac{w}{r}}, \quad L^* = \bar{Q}\sqrt{\frac{r}{w}}
\]

\textbf{Interpretation:} The firm allocates inputs such that the marginal product per dollar is equalized across inputs.

\section*{9. Applications in Economics}

\begin{center}
\begin{tabular}{l|l|l}
\textbf{Field} & \textbf{Objective Function} & \textbf{Constraint} \\ \hline
Consumer Theory & $U(x,y)$ & $p_xx + p_yy = M$ \\
Production Theory & $Q(K,L)$ & $wL + rK = C$ \\
International Trade & Welfare $W(t)$ & Balance of trade constraint \\
Public Finance & Utility & Government budget constraint \\
Growth Models & Output or Consumption & Capital accumulation equation
\end{tabular}
\end{center}




\end{document}
